\subsection{Qualitative Comparison}\label{subsec:qualitative_compare}

Figure~\ref{fig:window_sizes} presents the binarization results obtained with the different adaptive thresholding methods for a window size of $7\times 7$.
Overall, most approaches are able to separate foreground text from the noisy background, although notable differences in robustness and visual clarity can be observed.

% 7x7
For the smallest window size, Niblack's method produces the most distinct output. 
While the algorithm captures the textual structures reasonably well, it fails to suppress background noise effectively, leading to reduced contrast between the text in the image (foreground) and background.
Sauvola's method improves upon this behavior but still exhibits occasional inaccuracies in low-contrast regions.
In contrast, the Nick approach yields cleaner results despite being closely related to the Niblack and Sauvola formulations, indicating improved robustness of its modified threshold computation.
Among all methods Bradley's integral image based approach achieves the clearest separation ot text from background, producing the most visually consistent binarization. 

% 15x15
Increasing the window size to $15\times 15$ does not fundamentally change the relative performance of the methods, but it affects the amount of remaining background artifacts.
In particular, the modern approaches by Bradley and Bataineh tend to produce thicker character strokes and more pronounced contours.
Niblack's method benefits from the larger neighborhood, showing improved foreground-background separation compared to the $7\times 7$ case, although background noise remains more prominent than for the other approaches.

% 21x21 
For the largest evaluated window size of $21\times 21$, these effects become even more pronounced. 
While Niblack continues to improve in separating text from background, the output of Bataineh's method increasingly includes bold foreground structures as well as additional background contours.
This behavior is illustrated in Figure~\ref{fig:bataineh_window_sizes}, where larger windows lead to stronger stroke widths. 

Overall, Bradley's approach appears least sensitive to different window sizes, providing consistently high-quality binarization across all evaluated configurations.
Bataineh's method achieves strong contrast but requires careful parameter selection, whereas the classical methods show moderate improvements with increasing window size but remain less robust under heavy noise.
